\originalchapter{习题课十: 位移电流与能流}
题源为\href{https://ocw.mit.edu/courses/8-02-physics-ii-electricity-and-magnetism-spring-2007/resources/solving10/}{Problem Solving 10}, PDF第1--10页(原页脚沿用Solving 11). Q1--Q13和考试练习题(a)--(e)全部保留, 所有解答为AI独立解答. 原Q11把能流误回指Q2、Q12误回指Q3; 以下明确对应前一能流计算. Q13正文放电条件误印$\dot Q>0$, 答题页却为$\dot Q<0$, 按放电物理意义解释并保留勘误说明.

\begin{exercise}
充电电容器两圆板半径$a$, 间距$d\ll a$, 下板$+Q(t)$、上板$-Q(t)$, 正$z$向上. $I(t)=\dot Q(t)>0$, 忽略边缘, 场在$r>a$为零. 在两板中间取半径$r<a$、法向$+z$圆盘.
Q1求板间电场.
Q2求圆盘电通量.
Q3求圆盘位移电流.
Q4求真实传导电流.
Q5取圆盘边缘逆时针环路, 写磁场环量.
Q6求磁场大小.
Q7按电场方向伸右拇指, 磁场沿右指卷向还是反向?
Q8放电时磁场方向是否改变并解释.
Q9求$r\le a$的坡印廷矢量.
Q10在$r=a$、高$d$圆柱侧面求向外能流及单位.
Q11用$C=\eps\pi a^2/d$把前问改写为$Q,I,C$.
Q12核场能增加率等于侧面向内功率.
Q13若$Q>0$且在放电, 图像如何改变?
\end{exercise}
\begin{solution}
Q1. $\ee=Q\uv z/(\eps\pi a^2)$.
Q2. $\Phi_E=Q r^2/(\eps a^2)$.
Q3. $I_d=\eps\dot\Phi_E=I r^2/a^2$.
Q4. 真空中传导电流为零.
Q5. 轴对称磁场环向, 环量$2\pi rB_\phi$.
Q6. 安培--麦克斯韦式给$B_\phi=\mu_0Ir/(2\pi a^2)$. $I$为有符号电流; 此时$I>0$所以大小是该正值.
Q7. 与右手卷向相同, 从上方看逆时针.
Q8. 若仍$Q>0$但$I<0$, 磁场反向, 因为场由$\dot Q$而不是$Q$的正负决定.
Q9. $\uv z\times\uv\phi=-\uv r$, 所以
\[
\vct S=-\frac{QIr}{2\eps\pi^2a^4}\uv r,\qquad 0\le r<a.
\]
Q10. 以侧面向外为正, $P_{\rm out}=S_r(a)2\pi ad=-QId/(\eps\pi a^2)$, 单位W, 负值表示向内.
Q11. $P_{\rm out}=-QI/C$, 向内功率$QI/C$.
Q12. $U=Q^2/(2C)$, 对固定几何求导$\dot U=Q\dot Q/C=QI/C=-P_{\rm out}$, 无真空焦耳热时完全符合能量守恒.
Q13. 放电必须取$\dot Q=I<0$, 电场此刻仍向上, 磁场反向, $\vct S$向外, $U$下降. 原正文$\dot Q>0$与“$Q>0$放电”互相矛盾, 不能把它机械代入而宣称得到放电.
\end{solution}

\begin{exercise}
考试练习题. 同样圆板电容器初始下板$+Q_0$、上板$-Q_0$, 在$t=0$把半径$a$、高$d$的均匀圆柱电阻整块插入两板并良好接触, 电容通过内部电阻放电. 忽略边缘. 在板间取半径$r<a$、法向向上的圆盘.
(a)用高斯定律求电场并判上下方向.
(b)求圆盘向上传导电流, 用$Q$或$\dot Q$和所给参数表示.
(c)求电通量变化率.
(d)用安培--麦克斯韦定律求圆盘边缘磁场环量.
(e)解释该结果与能量耗散、能流是否相容.
\end{exercise}
\begin{solution}
(a)$\ee=Q(t)\uv z/(\eps\pi a^2)$向上.
(b)电阻中的正电流向上, 下板正电荷因此减少, 通过整块电阻的传导电流是$-\dot Q>0$. 均匀分布给$I_c(r)=-\dot Q\,r^2/a^2$.
(c)$\dot\Phi_E=\dot Q\,r^2/(\eps a^2)<0$, 所以$I_d=\dot Q\,r^2/a^2$.
(d)$\oint\bb\cdot\dd\vct\ell=\mu_0(I_c+I_d)=0$. 在无外加磁场且理想轴对称的模型中, 环向$B_\phi=0$.
(e)无矛盾. 与真空充电不同, 这里能量就在原有电场所在体积内转为热, 无须侧面传入. 设电导率$\sigma_{\rm c}$, 则$\dot Q=-\sigma_{\rm c}Q/\eps$, $\jj\cdot\ee=\sigma_{\rm c}E^2>0$, 而$\partial_tu_E=\eps E\dot E=-\sigma_{\rm c}E^2$. 场能减少恰好等于同处焦耳耗散, 坡印廷定理允许$\vct S=0$而仍有热. 若考虑插入的过渡过程、边缘或外部引线, 还要另计它们的局部场与能流.
\end{solution}

原答题页把圆柱侧面积误印为$2\pi ab$, 正确为$2\pi ad$（板间距记$d$）; 本题能流积分始终采用实际圆柱侧面积.
